\section{LİTERATÜR ARAŞTIRMASI VE BENZER SİSTEMLER}
\label{bolum3}

\subsection{Mevcut Eğitim Platformları}
\label{mevcut-e-itim-platformlar}

REST API güvenliği ve web uygulama güvenliği konusunda geliştirilen eğitim platformları, son yıllarda siber güvenlik eğitim süreçlerinde önemli bir yere sahip olmuştur. Özellikle uygulamalı öğrenme yaklaşımının yaygınlaşmasıyla birlikte kullanıcıların gerçek güvenlik açıklarını kontrollü laboratuvar ortamlarında deneyimleyebilmesine olanak sağlayan platformlar geliştirilmiştir. Bu platformlar sayesinde öğrenciler, geliştiriciler ve siber güvenlik araştırmacıları yalnızca teorik bilgi edinmekle kalmamakta; aynı zamanda saldırı tekniklerini, güvenlik açıklarının oluşma mantığını ve temel savunma yöntemlerini uygulamalı olarak inceleyebilmektedir.

Bu alandaki en yaygın ve öncü platformların başında OWASP Juice Shop, DVWA (Damn Vulnerable Web Application) ve PortSwigger Web Security Academy gelmektedir. Bu platformlar güvenlik eğitimi açısından önemli katkılar sağlamasına rağmen, modern REST API güvenliği konusunda bazı sınırlılıklara sahiptir.

OWASP tarafından geliştirilen OWASP Juice Shop, bilinçli olarak güvenlik açıkları içeren modern web uygulamalarından biridir. Platform, gerçek bir e-ticaret sistemi mantığıyla tasarlanmış olup kullanıcı giriş sistemi, ürün yönetimi, alışveriş işlemleri ve veri işleme gibi birçok gerçek dünya senaryosunu içermektedir. Sistem içerisinde authentication bypass, Broken Access Control, SQL Injection, Cross-Site Scripting (XSS), güvenlik yapılandırma hataları ve hassas veri sızıntıları gibi birçok farklı güvenlik açığı bulunmaktadır. Kullanıcılar bu açıkları çözerek sistemdeki “challenge” görevlerini tamamlamaktadır. Juice Shop’un en önemli avantajlarından biri modern web teknolojileri kullanılarak geliştirilmiş olması ve kullanıcıya gerçekçi saldırı senaryoları sunabilmesidir.

Bununla birlikte OWASP Juice Shop ağırlıklı olarak genel web uygulama güvenliği üzerine odaklanmaktadır. Platform içerisinde bazı API tabanlı senaryolar bulunsa da sistemin temel amacı doğrudan REST API güvenliğini karşılaştırmalı olarak analiz etmek değildir. Ayrıca platform içerisinde güvenli ve güvensiz API ortamlarının aynı anda sunulmaması, kullanıcıların saldırı ve savunma mekanizmalarını eş zamanlı olarak gözlemlemesini zorlaştırmaktadır. Bu nedenle kullanıcılar saldırının nasıl gerçekleştiğini görebilmekte; ancak aynı sistemin güvenli sürümünde savunma mekanizmalarının davranışı nasıl değiştirdiğini doğrudan analiz edememektedir.

Bir diğer önemli eğitim platformu olan DVWA (Damn Vulnerable Web Application), web güvenliği eğitiminde uzun yıllardır kullanılan klasik laboratuvar ortamlarından biridir. DVWA özellikle başlangıç seviyesindeki güvenlik araştırmacıları ve öğrenciler için SQL Injection, Reflected XSS, Stored XSS, Command Injection, File Inclusion ve CSRF gibi temel web güvenlik açıklarının öğrenilmesini amaçlamaktadır. Platform PHP tabanlı geliştirilmiş olup kullanıcıların güvenlik açıklarını uygulamalı şekilde test edebilmesine olanak sağlamaktadır. Ayrıca güvenlik seviyelerinin “low”, “medium” ve “high” olarak değiştirilebilmesi sayesinde kullanıcılar aynı güvenlik açığının farklı koruma seviyelerindeki davranışını gözlemleyebilmektedir.

Ancak DVWA daha çok geleneksel web uygulama güvenliği yaklaşımına odaklanmaktadır. Platformun geliştirildiği dönem itibarıyla REST API mimarileri günümüzdeki kadar yaygın olmadığı için modern API güvenlik problemlerini kapsamlı şekilde ele almamaktadır. Özellikle JWT tabanlı kimlik doğrulama sistemleri, API rate limiting mekanizmaları, token manipülasyonu, Broken Object Level Authorization (BOLA), modern access control problemleri ve API gateway mantığı gibi güncel REST API güvenlik konuları platform içerisinde yeterince yer almamaktadır. Bu nedenle DVWA, temel web güvenliği eğitimi açısından güçlü olsa da modern REST API güvenliği konusunda sınırlı kalmaktadır.

PortSwigger tarafından geliştirilen Web Security Academy ise günümüzde en kapsamlı uygulamalı web güvenliği eğitim platformlarından biri olarak kabul edilmektedir. Platform; SQL Injection, XSS, CSRF, SSRF, authentication bypass, OAuth güvenliği, erişim kontrolü problemleri ve modern API güvenliği gibi çok geniş bir konu yelpazesini kapsamaktadır. Kullanıcılar platform üzerinde interaktif laboratuvarlar çözerek gerçek saldırı senaryolarını deneyimleyebilmektedir. Özellikle Burp Suite entegrasyonu sayesinde HTTP istekleri üzerinde detaylı analiz yapılabilmesi, platformun profesyonel güvenlik eğitiminde yaygın olarak kullanılmasını sağlamaktadır.

Web Security Academy’nin REST API güvenliği konusunda sunduğu laboratuvarlar modern sistemlere daha yakın senaryolar içermektedir. Özellikle JWT manipülasyonu, API authentication problemleri, erişim kontrolü zafiyetleri ve güvenlik doğrulama eksiklikleri gibi konular platform içerisinde uygulamalı olarak gösterilmektedir. Ancak platform daha çok saldırı perspektifine odaklanmakta; aynı saldırıların güvenli sistem mimarisinde nasıl engellendiğini doğrudan karşılaştırmalı biçimde göstermemektedir. Kullanıcılar güvenlik açığını sömürebilmekte; ancak aynı sistemin güvenli versiyonunu inceleyerek savunma mekanizmalarının etkisini doğrudan analiz edememektedir.

Bu platformların ortak özelliklerinden biri, güvenlik açıklarının uygulamalı olarak öğrenilmesine önemli katkı sağlamalarıdır. Ancak mevcut sistemlerin büyük çoğunluğu ya yalnızca saldırı tarafına odaklanmakta ya da modern REST API güvenliğini tam kapsamlı biçimde ele almamaktadır. Ayrıca birçok platformda saldırı ve savunma mekanizmalarının aynı sistem üzerinde eş zamanlı karşılaştırılması mümkün değildir. Bu durum kullanıcıların güvenli backend geliştirme prensiplerini uygulamalı olarak anlamasını zorlaştırmaktadır.

Bu proje kapsamında geliştirilen laboratuvar sistemi ise mevcut platformlardan farklı olarak hem güvensiz hem de güvenli REST API ortamlarını aynı mimari üzerinde sunmayı hedeflemektedir. Böylece kullanıcılar aynı saldırının güvenlik önlemi olmayan bir sistemde nasıl başarılı olduğunu ve modern güvenlik mekanizmalarıyla güçlendirilmiş bir sistemde neden başarısız olduğunu doğrudan gözlemleyebilmektedir. Ayrıca proje özellikle REST API güvenliği üzerine odaklanarak JWT doğrulama, RBAC, WAF filtreleme, rate limiting ve modern API savunma mekanizmalarını uygulamalı biçimde göstermeyi amaçlamaktadır.

Sonuç olarak mevcut eğitim platformları siber güvenlik eğitimi açısından önemli katkılar sağlamakla birlikte; modern REST API güvenliği, saldırı-savunma karşılaştırması ve güvenli backend geliştirme farkındalığı konularında bazı eksikliklere sahiptir. Bu çalışma, söz konusu eksiklikleri azaltmayı ve REST API güvenliğini daha kapsamlı, karşılaştırmalı ve uygulamalı biçimde öğretmeyi hedefleyen bir laboratuvar ortamı sunmaktadır.

\subsection{Akademik Çalışmalar}
\label{akademik-al-malar}

REST API güvenliği ve uygulamalı siber güvenlik eğitimi konusunda yapılan akademik çalışmalar, modern web servislerinin güvenlik açıklarına karşı oldukça hassas olduğunu ve özellikle uygulamalı eğitim yöntemlerinin güvenlik farkındalığını artırmada önemli rol oynadığını göstermektedir. Literatürde yer alan araştırmalar; REST API sistemlerinde erişim kontrolü, kimlik doğrulama mekanizmaları, veri güvenliği ve saldırı tespit süreçleri gibi konuların günümüzde kritik önem taşıdığını ortaya koymaktadır.

OWASP kapsamında tanımlanan API güvenlik açıklarının gerçek sistemlerde yaygın olarak bulunduğu birçok akademik çalışmada doğrulanmıştır. Vargas ve ark. (2020) tarafından gerçekleştirilen çalışmada, yüksek trafik hacmine sahip REST API servisleri analiz edilmiş ve incelenen sistemlerin \%60’tan fazlasında yetersiz yetkilendirme (authorization) kontrolleri bulunduğu tespit edilmiştir. Özellikle Broken Object Level Authorization (BOLA) zafiyetlerinin, kullanıcıların yetkisi olmayan verilere erişebilmesine neden olduğu belirtilmiştir. Çalışmada geliştiricilerin büyük bölümünün kimlik doğrulama mekanizmasını uygulamasına rağmen nesne seviyesinde erişim kontrolünü yeterince sağlamadığı vurgulanmıştır. Bu durumun özellikle kullanıcı verileri, sipariş sistemleri ve müşteri kayıtları gibi hassas veri içeren API servislerinde ciddi güvenlik riskleri oluşturduğu ifade edilmiştir.

Benzer şekilde API güvenliği üzerine yapılan diğer araştırmalar, modern REST servislerinde güvenlik kontrollerinin çoğu zaman işlevsellik odaklı geliştirme süreçlerinin gerisinde kaldığını göstermektedir. Özellikle hızlı geliştirme süreçleri, mikro servis mimarilerinin yaygınlaşması ve bulut tabanlı sistemlerin artışı; güvenlik yapılandırmalarında hatalara yol açabilmektedir. Literatürde yer alan birçok çalışmada SQL Injection, JWT manipülasyonu, güvenlik doğrulama eksiklikleri, yetersiz rate limiting mekanizmaları ve hatalı erişim kontrolü problemlerinin modern API sistemlerinde en sık karşılaşılan güvenlik problemleri arasında olduğu belirtilmektedir.

Akademik çalışmalar yalnızca güvenlik açıklarını incelemekle kalmamış, aynı zamanda güvenlik eğitimi yöntemlerinin etkinliği üzerine de yoğunlaşmıştır. Jensen ve ark. (2019) tarafından gerçekleştirilen araştırmada, uygulamalı CTF (Capture The Flag) tabanlı eğitim ortamlarının siber güvenlik farkındalığını artırmada geleneksel teorik öğretim yöntemlerine kıyasla çok daha etkili olduğu ortaya konmuştur. Araştırmada öğrenciler iki farklı gruba ayrılmış; bir gruba klasik teorik güvenlik eğitimi verilirken diğer grup uygulamalı laboratuvar ortamlarında saldırı senaryoları çözmüştür. Sonuçlar, uygulamalı eğitim alan öğrencilerin güvenlik açıklarını tanıma, saldırı mantığını anlama ve savunma mekanizmalarını yorumlama konusunda daha yüksek başarı gösterdiğini ortaya koymuştur.

Bu çalışmalar, özellikle uygulamalı laboratuvar ortamlarının öğrenme sürecindeki önemini vurgulamaktadır. Teorik eğitim modellerinde öğrenciler genellikle güvenlik açıklarının yalnızca tanımlarını öğrenirken; uygulamalı laboratuvar ortamlarında gerçek saldırı senaryolarını deneyimleyebilmekte ve sistem davranışlarını doğrudan gözlemleyebilmektedir. Bu yaklaşımın, güvenlik farkındalığının kalıcılığı açısından daha etkili olduğu akademik olarak desteklenmektedir.

Literatürde Docker tabanlı laboratuvar ortamları üzerine yapılan çalışmalar da dikkat çekmektedir. Özellikle konteyner tabanlı sistemlerin siber güvenlik eğitiminde sağladığı avantajlar birçok araştırmada vurgulanmıştır. Docker teknolojisi sayesinde laboratuvar ortamlarının farklı işletim sistemlerinde aynı şekilde çalıştırılabilmesi, eğitim süreçlerinde standartlaşmayı kolaylaştırmaktadır. Ayrıca konteyner izolasyonu sayesinde güvenlik açıkları kontrollü biçimde çalıştırılabilmekte ve gerçek sistemlere zarar verme riski azaltılmaktadır.

Docker tabanlı izole laboratuvar ortamlarının, farklı donanım altyapılarında tekrar üretilebilir deney koşulları sağladığı ve bu özelliğin eğitim kalitesini önemli ölçüde artırdığı çeşitli akademik çalışmalarda belirtilmiştir. Özellikle tekrar üretilebilirlik (reproducibility) kavramı, siber güvenlik eğitiminde büyük önem taşımaktadır. Çünkü aynı laboratuvar ortamının farklı kullanıcılar tarafından aynı koşullarda çalıştırılabilmesi; eğitim süreçlerinin tutarlılığını ve güvenilirliğini artırmaktadır. Geleneksel sanal makine tabanlı laboratuvarlara kıyasla Docker konteynerlerinin daha hafif, hızlı ve taşınabilir olması da modern eğitim platformlarında yaygın şekilde tercih edilmesine neden olmaktadır.

Bunun yanında bazı akademik çalışmalar, güvenli ve güvensiz sistemlerin karşılaştırmalı olarak sunulmasının öğrenme sürecini olumlu etkilediğini göstermektedir. Kullanıcıların aynı saldırının hem savunmasız hem de korumalı sistem üzerindeki davranışını gözlemleyebilmesi, güvenlik mekanizmalarının çalışma mantığını daha iyi anlamalarını sağlamaktadır. Ancak literatürdeki mevcut platformların büyük çoğunluğu yalnızca saldırı odaklı laboratuvarlar sunmakta; güvenlik önlemlerinin sistem davranışını nasıl değiştirdiğini doğrudan göstermemektedir.

Bu çalışma, literatürde belirtilen eksiklikleri dikkate alarak tasarlanmıştır. Projede hem güvensiz hem de güvenli REST API ortamları aynı sistem içerisinde sunularak kullanıcıların saldırı ve savunma mekanizmalarını karşılaştırmalı biçimde analiz edebilmesi hedeflenmiştir. Ayrıca Docker tabanlı konteyner mimarisi kullanılarak tekrar üretilebilir, taşınabilir ve kontrollü bir laboratuvar ortamı oluşturulmuştur. Böylece proje, akademik çalışmalarda önerilen uygulamalı öğrenme yaklaşımını REST API güvenliği alanına uyarlayan eğitim odaklı bir sistem sunmaktadır.

Sonuç olarak literatürdeki çalışmalar; REST API güvenlik açıklarının modern sistemlerde oldukça yaygın olduğunu, uygulamalı laboratuvar ortamlarının güvenlik eğitimi açısından önemli avantajlar sağladığını ve Docker tabanlı tekrar üretilebilir sistemlerin eğitim kalitesini artırdığını göstermektedir. Bu proje ise söz konusu akademik bulgular doğrultusunda geliştirilmiş, REST API güvenliğini uygulamalı ve karşılaştırmalı biçimde öğretmeyi amaçlayan bir laboratuvar ortamı sunmaktadır.

\subsection{API Fortress'in Farklılaştığı Noktalar}
\label{api-fortress-in-farkl-la-t-noktalar}

Bu çalışma kapsamında geliştirilen API Fortress sistemi, mevcut web güvenliği ve siber güvenlik eğitim platformlarından farklı olarak özellikle REST API güvenliği üzerine odaklanan, saldırı ve savunma mekanizmalarını aynı sistem içerisinde karşılaştırmalı biçimde sunan ve Docker tabanlı taşınabilir laboratuvar mimarisi kullanan uygulamalı bir eğitim platformu olarak tasarlanmıştır. Literatürde yer alan birçok platform güvenlik açıklarının sömürülmesine odaklanırken, API Fortress yalnızca saldırı mantığını öğretmeyi değil; aynı zamanda modern backend güvenlik mekanizmalarının sistem davranışını nasıl değiştirdiğini uygulamalı olarak göstermeyi hedeflemektedir.

API Fortress’in mevcut platformlardan ayrıldığı en önemli noktalardan biri, doğrudan REST API güvenliği üzerine yoğunlaşmasıdır. Günümüzde modern yazılım sistemlerinin büyük çoğunluğu frontend-backend ayrımına dayalı mimariler kullanmakta ve sistemler arasındaki veri iletişimi REST API servisleri üzerinden gerçekleştirilmektedir. Mobil uygulamalar, mikro servis mimarileri, bulut tabanlı sistemler ve modern web uygulamaları yoğun biçimde API teknolojilerine bağımlı hale gelmiştir. Buna rağmen mevcut eğitim platformlarının büyük bölümü hâlâ geleneksel web güvenliği açıklarına odaklanmakta; modern REST API zafiyetlerini kapsamlı şekilde ele almamaktadır.

API Fortress ise doğrudan OWASP API Security Top 10 kapsamındaki modern API güvenlik problemlerini hedef almaktadır. Sistem içerisinde Broken Object Level Authorization (BOLA), Broken Authentication, Security Misconfiguration, Injection saldırıları, yetersiz erişim kontrolü, token manipülasyonu, rate limiting eksiklikleri ve güvenli doğrulama problemleri gibi güncel REST API zafiyetleri uygulamalı olarak sunulmaktadır. Böylece kullanıcılar yalnızca klasik web uygulama güvenliğini değil; modern backend mimarilerinde karşılaşılan API güvenlik problemlerini de gerçekçi senaryolar üzerinden deneyimleyebilmektedir.

API Fortress’in farklılaştığı ikinci temel nokta, saldırı ve savunma ortamlarını aynı sistem içerisinde birlikte sunmasıdır. Mevcut güvenlik laboratuvarlarının büyük çoğunluğu yalnızca saldırı perspektifine odaklanmaktadır. Kullanıcılar güvenlik açıklarını sömürmeyi öğrenebilmekte; ancak aynı sistemin güvenli sürümünde savunma mekanizmalarının nasıl çalıştığını doğrudan gözlemleyememektedir. Bu durum, güvenli yazılım geliştirme süreçlerinin yeterince anlaşılmasını zorlaştırmaktadır.

API Fortress içerisinde ise aynı işlevleri yerine getiren iki farklı backend ortamı bulunmaktadır. İlk ortam kasıtlı olarak güvenlik açıkları içeren “vulnerable API” yapısını temsil ederken; ikinci ortam modern güvenlik mekanizmalarıyla güçlendirilmiş “secure API” yapısını temsil etmektedir. Bu yaklaşım sayesinde kullanıcılar aynı saldırının iki farklı sistem üzerindeki davranışını karşılaştırmalı olarak inceleyebilmektedir.

Örneğin kullanıcılar güvensiz ortamda SQL Injection saldırısının başarılı şekilde çalıştığını gözlemleyebilirken; güvenli ortamda input validation, filtreleme mekanizmaları ve güvenli sorgu yapıları nedeniyle aynı saldırının neden başarısız olduğunu analiz edebilmektedir. Benzer şekilde JWT manipülasyonu, Broken Authentication veya yetkisiz veri erişimi gibi saldırıların güvenli backend mimarisinde nasıl engellendiği doğrudan incelenebilmektedir. Bu karşılaştırmalı yaklaşım, kullanıcıların yalnızca saldırı mantığını değil; güvenlik önlemlerinin sistem davranışı üzerindeki etkisini de anlamasını sağlamaktadır.

Sistem aynı zamanda modern savunma mekanizmalarını uygulamalı olarak göstermektedir. Güvenli API ortamında JWT tabanlı kimlik doğrulama sistemi, RBAC (Role-Based Access Control), regex tabanlı temel WAF filtreleme sistemi, rate limiting mekanizması, güvenli input validation işlemleri ve CORS güvenlik yapılandırmaları uygulanmıştır. Böylece kullanıcılar teorik güvenlik prensiplerini gerçek backend mimarisi üzerinde inceleyebilmekte ve modern güvenli yazılım geliştirme yaklaşımını uygulamalı biçimde deneyimleyebilmektedir.

API Fortress’in üçüncü önemli farklılığı ise Docker tabanlı taşınabilir laboratuvar mimarisidir. Geleneksel güvenlik laboratuvarlarında sistem kurulumu çoğu zaman karmaşık olmakta; işletim sistemi bağımlılıkları, paket sürüm uyuşmazlıkları ve yapılandırma problemleri eğitim süreçlerini zorlaştırabilmektedir. Özellikle farklı donanım ve yazılım altyapılarında aynı laboratuvar ortamının tekrar oluşturulması önemli bir problem oluşturmaktadır.

API Fortress bu problemi Docker konteyner teknolojisi kullanarak çözmektedir. Sistem içerisindeki güvenli ve güvensiz backend servisleri konteyner tabanlı çalıştırılmakta ve tüm bağımlılıklar izole şekilde paketlenmektedir. Böylece kullanıcılar sistemi farklı işletim sistemlerinde veya geliştirme ortamlarında minimum kurulum işlemiyle çalıştırabilmektedir. Docker mimarisi aynı zamanda laboratuvar ortamının tekrar üretilebilirliğini artırmakta ve eğitim süreçlerinde standartlaşma sağlamaktadır.

Konteyner mimarisinin bir diğer avantajı ise güvenlik izolasyonudur. Güvensiz API servisleri kontrollü konteyner ortamlarında çalıştırıldığı için kullanıcıların saldırı testleri gerçek sistemlere zarar vermeden gerçekleştirilebilmektedir. Bu durum hem eğitim güvenliğini artırmakta hem de saldırı senaryolarının kontrollü şekilde uygulanabilmesine olanak sağlamaktadır.

API Fortress ayrıca eğitim erişilebilirliği açısından da önemli avantajlar sunmaktadır. Sistem herhangi bir bulut ortamına, lokal geliştirme ortamına veya sanal laboratuvar altyapısına kolayca taşınabilmektedir. Böylece öğrenciler, araştırmacılar ve geliştiriciler farklı cihazlar üzerinde aynı laboratuvar ortamını çalıştırabilmekte ve standart bir eğitim deneyimi elde edebilmektedir.

Mevcut platformların önemli bir kısmı yalnızca challenge çözmeye odaklanırken API Fortress aynı zamanda backend mimarisinin anlaşılmasını da hedeflemektedir. Kullanıcılar yalnızca saldırıyı gerçekleştirmekle kalmamakta; API endpoint yapıları, JWT doğrulama süreçleri, rate limiting mantığı, WAF filtreleme mekanizmaları ve erişim kontrol sistemleri gibi backend bileşenlerinin çalışma mantığını da inceleyebilmektedir. Bu yönüyle sistem yalnızca bir saldırı laboratuvarı değil; aynı zamanda modern backend güvenliği eğitim platformu niteliği taşımaktadır.

Sonuç olarak API Fortress; doğrudan REST API güvenliği üzerine odaklanması, saldırı ve savunma ortamlarını karşılaştırmalı biçimde sunması, modern güvenlik mekanizmalarını uygulamalı olarak göstermesi ve Docker tabanlı taşınabilir laboratuvar mimarisi kullanması sayesinde mevcut eğitim platformlarından ayrışmaktadır. Bu özellikler sayesinde sistem, REST API güvenliği konusunda hem teknik hem de eğitimsel açıdan kapsamlı ve uygulama odaklı bir öğrenme ortamı sunmaktadır.


\clearpage
